% !TEX root = thesis.tex

\chapter{Preliminaries and Taxonomy}\label{chap:preliminaries-and-taxonomy}

\section{Transformer Architecture}

The Transformer model is one of the most groundbreaking technologies in the field of Natural Language Processing (NLP). Unlike sequential traditional models such as RNNs (Recurrent Neural Networks), the Transformer architecture uses an innovative attention mechanism to keep track of long-range dependencies between inputs and outputs. What makes this possible is a clever mechanism, which holds contextual meaning during processing of the input: Multi-Head Self-Attention  ~\cite{vaswani2017attention}

The Multi-Head Self-Attention actually builds up on something more fundamental: the Sefl-Attention mechanism. For a given set of input tokens, the Self-Attention mechanism transforms the initial token embeddings into contextualised representations, where the representation of each token incorporates information from other relevant tokens in the input sequence. To calculate self-attention, Key, Query and Value vektors generated by multiplying input tokens with learned weights are used as shown in the following formula: 
\[ 
  Attention(Q,K,V) = softmax\bigg(\frac{QK^T}{\sqrt{d_k}}\bigg)
\]
where, Q, K, and V represent the query, key, and value matrices, respectively, and $d_k$ is the dimension of the key vectors.

The Multi-Head Self-Attention as afore mentioned, builds on this self-attention mechanism, by carrying out the same operation in parallel with different learned weights and concatenating the output of each individual head, in order to capture other possible contextual meanings of input tokens which would not necessarily be captured by a single attention head. This can be represented as follows:

\[ 
  MultiHead(Q,K,V) = Concat(head_1,head_2,...,head_h) W^O
\]

where \( head_i = Attention(QW_i^Q, KW_i^K,VW_i^V) \) stands for the \textit{i}-th attention head, and $W^O$ is a learned linear transformation \cite{bai2024beyond} 

As it can be seen from the above equations Attention presents one of the bottlenecks of efficient LLM inference, as it is a quite compute intensive operation.


\section{Inference in LLMs}

The inference pipeline in Transformer decoder based LLMs is made up of mainly two stages: the Prefill Stage and Decode Stage  

The prefill stage consists of creating Key-Value (KV) pairs from the input tokens, which will later be cached (KV-cache). These are used to determine the contextual importance of one token when compared with another. At this point it is worth highlighting that the size of the KV-cache grows with input length, which could easily strain memory. \cite{zhou2024efficientinf}

The Decode Stage consists of computing new Key and Value pairs based on the output tokens generated  and appending these to the alreading existing KV-cache. \cite{yuan2024llminfunv}

A brief overview of the how inference works in LLMs, can be seen in Fig. 2.1


\begin{figure}[ht]
\centering
\begin{tikzpicture}[
  node distance   = 0.75cm and 0.5cm,
  process/.style  = {rectangle, draw=blue!70, fill=blue!8,
                     text width=6.8cm, text centered,
                     rounded corners=3pt, minimum height=1.5em, font=\small},
  highlight/.style= {rectangle, draw=red!70, fill=red!8,
                     text width=6.8cm, text centered,
                     rounded corners=3pt, minimum height=1.5em, font=\small},
  io/.style       = {rectangle, draw=green!60!black, fill=green!8,
                     text width=6.8cm, text centered,
                     rounded corners=3pt, minimum height=1.5em, font=\small},
  annot/.style    = {rectangle, draw=red!50, fill=red!5,
                     text width=4.6cm, align=left,
                     rounded corners=2pt, font=\footnotesize},
  decision/.style = {diamond, draw=orange!80, fill=orange!8,
                     text width=2.0cm, text centered,
                     aspect=2.8, font=\small},
  arrow/.style    = {-{Stealth}, thick},
  darrow/.style   = {dashed, -{Stealth}, thick, red!50}
]
%% --- Main flow ---
\node [io]       (inp)    {Input Prompt (text)};
\node [process,  below=of inp]    (tok)    {(1)\ Tokenisation};
\node [process,  below=of tok]    (emb)    {(2)\ Token Embedding + Positional Encoding};
\node [highlight, below=of emb]   (attn)   {(3)\ Multi-Head Self-Attention + \ac{kv} Cache};
\node [highlight, below=of attn]  (ffn)    {(4)\ Feed-Forward Network (FFN)};
\node [process,  below=of ffn]    (lmh)    {(5)\ LM Head: Linear + Softmax};
\node [decision, below=of lmh, node distance=1.1cm] (eos)  {EOS?};
\node [io,       below=of eos, node distance=1.1cm]  (out)   {Generated Output (text)};

%% --- Arrows main flow ---
\foreach \from/\to in {inp/tok, tok/emb, emb/attn, attn/ffn, ffn/lmh, lmh/eos}
  \draw [arrow] (\from) -- (\to);
\draw [arrow] (eos.south) -- node[right, font=\footnotesize]{Yes} (out.north);
\draw [arrow] (eos.east)  -- ++(2.2,0) |- node[pos=0.25, right, font=\footnotesize]{No (decode loop)} (attn.east);

%% --- Bottleneck annotations ---
\node [annot, right=1.8cm of attn] (ba) {%
  \textbf{Bottleneck:}\\
  $\mathcal{O}(n^2)$ memory \& compute\\
  \ac{kv} cache $\propto L \!\cdot\! n$\\
  Bandwidth-bound in decode};
\node [annot, right=1.8cm of ffn]  (bf) {%
  \textbf{Bottleneck:}\\
  ${\sim}2/3$ of all parameters\\
  Bandwidth-bound in decode};

\draw [darrow] (attn.east) -- (ba.west);
\draw [darrow] (ffn.east)  -- (bf.west);

%% --- Phase labels ---
\node [left=1.4cm of emb, font=\footnotesize\color{blue!70}, align=right,
       text width=2.2cm] {\textit{Prefill phase}\\(parallel)};
\node [left=1.4cm of attn, font=\footnotesize\color{red!70}, align=right,
       text width=2.2cm] {\textit{Decode phase}\\(sequential)};
\end{tikzpicture}
\caption{Schematic of the Transformer inference pipeline. Red-shaded steps (3, 4)
  represent the primary resource bottlenecks. The decode loop repeats steps 3--5
  until an end-of-sequence token is produced.}
\label{fig:inference_pipeline}
\end{figure}


\section{Taxonomy and Key Metrics}\label{sec:taxonomy}

The following terms appear consistently throughout this paper.

\begin{description}
  \item[Computation:] Amount of processing power used usually measured in Floating Point Operations (FLOPs).
  \item[Memory:] Amount of storage space or RAM.
  \item[Energy:] Amount of electrical power used up.
  \item[Latency:] Time it takes on average to generate a token
  \item[Bandwidth:] How quickly data moves from one place to another. In our case, this could for example be how fast the parameters of the LLM are loaded into the GPU.
  \item[\ac{kv} cache:] The stored key and value tensors for all past tokens, maintained to avoid recomputation at each decode step.
  \item[Quantisation:] Representing weights and/or activations at reduced numerical precision to reduce memory footprint and, on hardware with native low-precision arithmetic, increase throughput.
  \item[Pruning:] Removing parameters from a trained model. \emph{Unstructured} pruning zeros individual weights; \emph{structured} pruning removes entire rows, columns, heads, or layers to yield a dense, smaller model directly executable on standard hardware.
  \item[Knowledge Distillation:] Training a compact \emph{student} model to match the output distribution of a larger \emph{teacher}, transferring capability at a fraction of the inference cost.
  \item[Throughput:] Tokens (or requests) completed per second across all concurrent users; the primary efficiency metric for server-side deployment.
  \item[Latency:] Elapsed time from request arrival to final-token delivery. \ac{ttft} (time to first token) governs perceived responsiveness; \ac{tpot} (time per output token) governs generation speed.
  \item[Arithmetic intensity:] Ratio of floating-point operations to bytes of memory traffic, in OPs/byte. Kernels below the hardware ridge point are bandwidth-bound; those above are compute-bound.
  \item[Edge device:] Any computing platform with constrained memory, power, and compute operating outside a data centre: smartphones or embedded controllers.
\end{description}

